% !TeX root = ../../Book1.tex
% 纲要标题：### 22.1 无限代数扩张
% 纲要位置：计划纲要.md，第 769 行。

\section{无限代数扩张}
\label{b1:sec:2201}

无限代数扩张中的每一个元素，仍落在某个有限子扩张中。困难不在单个元素，而在于不同有限层上的选择必须彼此相容。本节先把自同构写成一族有限层上的自同构，再说明为什么任意选择各层自同构通常不能拼成一个整体。

% 纲要内容（仅作写作计划，尚非教材正文）：
%
% * 有限 Galois 子扩张的有向系统；说明有限层须被全局自同构保持，并区分过渡映射与整体限制映射的满射性
% * 代数闭包的自同构群；解释可分闭包与纯不可分部分，以及自同构群同构为何不能代替扩张的可分性
% * 逐层相容的自同构；证明拼合与所选有限层无关
%

\subsection{有限子扩张的有向系统}
\label{b1:subsec:2201:01}
本章固定代数扩张 \(L/K\)。若它正规且可分，仍称 \(L/K\) 为 Galois 扩张，但不再要求次数有限；这时把它的自同构群记为 \(G=\operatorname{Gal}(L/K)\)。有限层上的群阶仍等于相应的扩张次数，整个无限扩张的域次数与自同构群的普通基数却一般不相等。记录各个有限 Galois 层上的自同构及其限制映射，能够保留有限理论中可以逐层计算的信息；这些有限层还将决定群上的拓扑。

\begin{proposition}\label{b1:prop:finite-galois-layers}
设 \(L/K\) 为代数 Galois 扩张。\(L\) 的每个有限子集都包含在某个有限 Galois 子扩张 \(N/K\) 中。所有这样的 \(N\) 按包含关系构成有向集合，且 \(L=\bigcup_N N\)。
\end{proposition}
\begin{proof}
取有限集 \(a_1,\ldots,a_r\in L\)。每个 \(a_i\) 的最小多项式可分，且由正规性在 \(L\) 中分裂。只有有限多个这样的多项式，每个又只有有限多个根，所以它们的全部根组成有限集。将这些根加入 \(K\)，所得域 \(N\) 仍在 \(L\) 内，由有限多个代数元生成，故 \(N/K\) 有限。把相同的首一最小多项式只保留一次，所得乘积可分，因为不同不可约多项式没有公共根；\(N\) 正是该乘积的分裂域，因此正规且可分，并含原有限集。空集可取 \(N=K\)。

若 \(N_1,N_2\) 是两个有限 Galois 子扩张，各取有限生成元，将它们的并按上一段处理，就有同时包含两域的有限 Galois 子扩张。这就是按包含关系有向的含义；有限次重复，任意有限多个这样的域都能放进一个共同有限 Galois 层中。每个 \(a\in L\) 所成的单点集也被某层包含，故它们的并为 \(L\)。
\end{proof}
有向集合的抽象定义见\cref{b1:def:directed-system}；这里的共同上层已由上述证明构造。记这个有向集合为 \(\mathcal N\)，并记 \(G_N=\operatorname{Gal}(N/K)\)。选取有限层时要求 \(N/K\) 正规，是为了能在它上面记录整体自同构的限制：由\cref{b1:thm:normal-embedding-criterion}，对任意 \(\sigma\in G\)，有 \(\sigma(N)=N\)。若只取任意有限子扩张 \(E/K\)，\(\sigma(E)\) 可能是另一个共轭子域，\(\sigma|_E\) 就不能视为 \(E\) 的自同构。同理，对 \(N\subseteq M\)，\(M\) 的每个 \(K\)-自同构都保持 \(N\)，所以限制给出群同态
\[
\operatorname{res}_{M,N}:G_M\longrightarrow G_N.
\]
这个有限层间的限制映射是否满射，以及 \(G\) 能否实现有限层上的每个自同构，都还需要嵌入延拓定理保证。

\begin{lemma}\label{b1:lem:galois-restriction-surjective}
设 \(L/K\) 为代数 Galois 扩张，\(G=\operatorname{Gal}(L/K)\)。若 \(N\subseteq M\subseteq L\) 都是 \(L/K\) 的有限 Galois 子扩张，并记 \(G_N=\operatorname{Gal}(N/K)\)、\(G_M=\operatorname{Gal}(M/K)\)，则限制同态
\[
G\longrightarrow G_N,
\qquad
\operatorname{res}_{M,N}:G_M\longrightarrow G_N
\]
都满射。
\end{lemma}
\begin{proof}
取包含 \(L\) 的 \(K\) 的代数闭包 \(\Omega\)。给定 \(\tau\in G_N\)，把它看作 \(N\rightarrow\Omega\) 的嵌入。\cref{b1:thm:embedding-extension}应用于代数扩张 \(L/N\)，将它延拓为 \(\sigma:L\rightarrow\Omega\)。因为 \(\sigma|_N=\tau\)，它固定 \(K\)；\cref{b1:thm:normal-embedding-criterion}及 \(L/K\) 的正规性给出 \(\sigma(L)=L\)，故 \(\sigma\in G\)，且它在 \(N\) 上恰为 \(\tau\)。这证明 \(G\rightarrow G_N\) 满射，即有限层的每个自同构都能在整体扩张上实现。

把同一论证中的 \(L\) 换成 \(M\)，\(M/N\) 仍代数而 \(M/K\) 正规，就得到 \(\tau\) 到 \(M\) 的自同构延拓。因此 \(G_M\rightarrow G_N\) 也满射，即较小有限层上的自同构能提升到较大有限 Galois 层。第一种满射控制整体群在每个有限坐标上的像，第二种控制有限层之间的过渡映射。对于三层域，逐次限制等于直接限制。
\end{proof}
这两种满射性来自延拓与正规性，并不是把各有限群写成坐标以后自动得到的性质。

\subsection{代数闭包的自同构群}
\label{b1:subsec:2201:02}
不完美域的代数闭包中，可能有基域之外的元素被每个基域自同构固定。例如 \(K=\mathbb F_p(t)\)，在代数闭包中取 \(u=t^{1/p}\)。多项式 \(X^p-t\) 在 \(K\) 上不可约，且只有 \(u\) 这一个不同根，所以 \(u\notin K\)，每个 \(K\)-自同构却都固定 \(u\)。因此 \(\operatorname{Aut}_K(\overline K)\) 在整个代数闭包上的不动域可以严格大于 \(K\)。要使自同构的作用能够检测扩张中的不同共轭根，需要考察可分部分。

\ProseParagraphBreak
固定 \(K\) 的代数闭包 \(\overline K\)。其中所有在 \(K\) 上可分的元素组成子域，记为 \(K^{\mathrm{sep}}\)，称为 \(K\) 的\term[kefen-bibao]{可分闭包}[separable closure]。它是子域的依据为\cref{b1:prop:separable-subfield}；任一可分元的全部共轭根仍然可分，且已在 \(\overline K\) 中，所以 \(K^{\mathrm{sep}}/K\) 正规而可分。
\symidx{\(K^{\mathrm{sep}}\)：所选代数闭包中的可分闭包}

\begin{definition}\label{b1:def:absolute-galois-group}
设 \(K\) 为域，固定一个代数闭包 \(\overline K\)，并令 \(K^{\mathrm{sep}}\subseteq\overline K\) 为所有在 \(K\) 上可分的元素组成的可分闭包。群
\[
G_K\coloneqq\operatorname{Gal}(K^{\mathrm{sep}}/K)
\]
称为域 \(K\) 的\term[juedui-galois-qun]{绝对 Galois 群}[absolute Galois group]。换一份可分闭包会给出同构的群，但同构一般依赖于闭包之间同构的选择。
\end{definition}

\symidx{\(G_K\)：域的绝对 Galois 群}

\begin{proposition}\label{b1:prop:closure-automorphisms}
设 \(K\) 为域，\(\overline K\) 为 \(K\) 的代数闭包，\(K^{\mathrm{sep}}\subseteq\overline K\) 为其中所有在 \(K\) 上可分的元素组成的子域，并令 \(G_K=\operatorname{Gal}(K^{\mathrm{sep}}/K)\)。则限制映射 \(\operatorname{Aut}_K(\overline K)\rightarrow G_K\) 为群同构。若 \(K\) 完美，则 \(K^{\mathrm{sep}}=\overline K\)；一般情形中不能把 \(\overline K/K\) 直接视为 Galois 扩张。
\end{proposition}
\begin{proof}
\(K\)-自同构保持最小多项式及其可分性，故保持可分闭包，限制映射有定义。若特征为零，所有代数元可分，限制映射就是恒等识别。以下设特征为 \(p>0\)。对任意 \(a\in\overline K\)，\cref{b1:prop:inseparable-polynomial-shape}把它的最小多项式写为
\[
m_{a,K}(X)=g(X^{p^r}),
\]
其中 \(g\) 首一、不可约且可分。代入 \(a\) 得 \(g(a^{p^r})=0\)，故 \(g\) 正是 \(a^{p^r}\) 的最小多项式。因此 \(a^{p^r}\) 在 \(K\) 上可分，属于 \(K^{\mathrm{sep}}\)。这说明每个 \(a\) 都在 \(K^{\mathrm{sep}}\) 上纯不可分。

若两个 \(K\)-自同构 \(\sigma,\tau\) 在 \(K^{\mathrm{sep}}\) 上有相同限制 \(\nu\)，对给定 \(a\) 取上述的 \(r\)，则
\[
\sigma(a)^{p^r}=\nu(a^{p^r})=\tau(a)^{p^r}.
\]
两像都是方程 \(X^{p^r}-\nu(a^{p^r})=0\) 的根。这个方程至多有一个不同根，因为两根之差的 \(p^r\) 次幂为零，而域没有非零幂零元。因此 \(\sigma(a)=\tau(a)\)；对每个 \(a\) 都如此，限制映射便单射。

反过来，给定 \(\nu\in G_K\)，把它看作 \(K^{\mathrm{sep}}\rightarrow\overline K\) 的嵌入。\cref{b1:thm:embedding-extension}应用于代数扩张 \(\overline K/K^{\mathrm{sep}}\)，将它延拓为 \(\sigma:\overline K\rightarrow\overline K\)。这个嵌入固定 \(K\)，代数闭包在 \(K\) 上正规，\cref{b1:thm:normal-embedding-criterion}保证 \(\sigma(\overline K)=\overline K\)，所以它是自同构，得到满射性。完美域的所有代数元可分，最后的等式由定义得到。
\end{proof}
纯不可分部分上的延拓唯一，因而没有增加自同构的选择，这正是两个群同构的原因。但这个同构不会使原本不可分的元素变得可分；这些不能被移动的元素仍在代数闭包中，使该群在那里作用的不动域可以严格大于 \(K\)。所以，还须区分群本身与它所作用的域，不能由上述群同构把 \(\overline K/K\) 直接当作 Galois 扩张。

\subsection{相容自同构}
\label{b1:subsec:2201:03}
从有限层上的自同构恢复整体映射，自然要对 \(a\in L\) 规定 \(\sigma(a)=\sigma_N(a)\)，其中 \(N\) 是任意包含 \(a\) 的有限 Galois 层。首先须解决的是选择无关性：同一个元素可能属于许多层，只有这些层给出的像一致，才能在 \(L\) 上定义一个映射。

\begin{proposition}\label{b1:prop:galois-compatible-family}
设 \(L/K\) 为代数 Galois 扩张，\(G=\operatorname{Gal}(L/K)\)，并令 \(\mathcal N\) 为 \(L/K\) 的全部有限 Galois 子扩张组成的有向集合。对每个 \(N\in\mathcal N\)，记 \(G_N=\operatorname{Gal}(N/K)\)，并给定 \(\sigma_N\in G_N\)。则存在唯一 \(\sigma\in G\) 在每个 \(N\) 上等于 \(\sigma_N\)，当且仅当
\begin{equation}\label{b1:eq:galois-compatibility}
\sigma_M|_N=\sigma_N\qquad(N\subseteq M).
\end{equation}
\end{proposition}
\begin{proof}
限制一个自同构当然满足相容条件。反过来，对 \(a\in L\)，选含 \(a\) 的有限 Galois 层 \(N\)，规定 \(\sigma(a)=\sigma_N(a)\)。若还选了 \(M\)，由\cref{b1:prop:finite-galois-layers}取同时包含两者的一层 \(P\)，相容性给出 \(\sigma_N(a)=\sigma_P(a)=\sigma_M(a)\)，所以定义无关选择。

对任意两个元素，取包含它们的同一有限层，就可在该层验证保持和、积与单位，以及固定 \(K\)。还须保证所得映射可逆。若 \(N\subseteq M\)，相容性给出 \(\sigma_M|_N=\sigma_N\)，而 \(\sigma_N\) 是 \(N\) 的自同构；所以 \(\sigma_M^{-1}\) 在 \(N\) 上的限制正是 \(\sigma_N^{-1}\)。因此逆族 \((\sigma_N^{-1})_N\) 也相容，同样拼成一个映射，在每层上都是 \(\sigma\) 的双边逆，因而在 \(L\) 上也是双边逆。所得 \(\sigma\) 是自同构；\(L\) 是各层的并，全部限制也就唯一确定了它。
\end{proof}
在每层能够延拓，不意味着可以把所有层上任意选定的自同构拼在一起。例如在有限域的塔 \(\mathbb F_q\subset\mathbb F_{q^2}\subset\mathbb F_{q^4}\) 中，取四次层上的 Frobenius \(F(x)=x^q\)。它在二次层上的限制仍为 \(x\mapsto x^q\)；由\cref{b1:thm:finite-field-automorphisms}，这个限制的阶为 \(2\)，是二次扩张的非平凡自同构。所以若预先在二次层选择恒等，它就不能与四次层上选定的 \(F\) 相容，违反\eqref{b1:eq:galois-compatibility}。各层的选择分别合法，还须满足限制相容条件才能拼成整体。下一节把相容族组成的群称为逆极限，并在其上建立适合这些有限信息的拓扑。
